\documentclass[11pt]{article}
\usepackage{fontspec}
\setmainfont{Times New Roman}
\setsansfont{Arial}
\setmonofont{Consolas}
\usepackage{amsmath,amssymb}
\usepackage{geometry}
\usepackage{microtype}
\geometry{a4paper,margin=3cm}
\setlength{\parindent}{1.25em}
\setlength{\parskip}{0pt}
\emergencystretch=10em
\allowdisplaybreaks
\providecommand{\Kfield}{\mathfrak{K}}
\providecommand{\Ssys}{\mathfrak{S}}
\providecommand{\Jdom}{\mathfrak{J}}
\providecommand{\Lsys}{\mathfrak{L}}
\providecommand{\Mfield}{\mathfrak{M}}
\providecommand{\tq}{\tau}
\providecommand{\ds}{\displaystyle}

\begin{document}

\section*{\S\ 7. Libovolny sistem $\Ssys_{n\rho}$, linearna familija $\Lsys_{n\rho}$ i integralna oblast $\Jdom_{n\rho}$ racionalnyh funkcij. Eksistencija racionalnoj bazy.}

Iz eksistencije racionalnoj bazy polj $\Kfield_{n\rho}$ teper bude vyvedena racionalna baza za libovolne sistemy racionalnyh i cělyh racionalnyh funkcij. My uporędžujemo te sistemy tako, kako one sukcesivno nastavaju jedna iz drugoj črez elementarne operacije dodavanja, množenja i děljenja.

I. Nehaj, kako vsegda, $\Ssys_{n\rho}$ označuje libovolny sistem racionalnyh (ili cělyh racionalnyh) funkcij od $n$ neopredělennyh, algebraičnogo ranga $\rho$. Elementy iz $\Omega$ takože računajut se kako prinadležeče sistemu.

II. Sistem $\Lsys_{n\rho}$ zove se {\itshape linearna familija}, ako ispolnja uslovja:

1. skupaj s $f(x)$ sistem $\Lsys_{n\rho}$ vsegda soderži takože $c\cdot f(x)$, kde $c$ jest libovolny element iz $\Omega$;

2. skupaj s $f(x)$ i $g(x)$ sistem $\Lsys_{n\rho}$ vsegda soderži takože $f(x)+g(x)$.

III. Linearna familija $\Jdom_{n\rho}$ zove se {\itshape integralna oblast} pod daljnym uslovjem:

3. skupaj s $f(x)$ i $g(x)$ oblast $\Jdom_{n\rho}$ vsegda soderži takože $f(x)\cdot g(x)$.

IV. Integralna oblast $\Kfield_{n\rho}$ zove se {\itshape polje} pod daljnym uslovjem:

4. zajedno s $f(x)$ i $g(x)$ --- kde $g(x)\ne0$ --- polje $\Kfield_{n\rho}$ vsegda soderži takože kvocient $f(x):g(x)$.

Uslovja 1--4 sut ty same, ktore v \S\ 1 byly dane za polje.\footnote{$f(x)$ i $g(x)$ mogut byti takože nultogo stepena, to jest elementy iz $\Omega$.}

Te oblasti treba teper sukcesivno vyvesti jedna iz drugoj.

a) Libovolny sistem $\Ssys_{n\rho}$ razširja se do linearnoj familije $\Lsys_{n\rho}$ slědujućim potrěbovanjem:

$\Lsys_{n\rho}$ soderži, kromě $\Ssys_{n\rho}$, vse funkcije $h(x)$ i samo te, ktore možno linearno predstaviti s koeficientami iz $\Omega$ črez konečno čislo funkcij iz $\Ssys_{n\rho}$:
\[
  h(x)=c_1f_1(x)+c_2f_2(x)+\cdots+c_kf_k(x),
\]
pričem $f_1(x),\ldots,f_k(x)$ proběgajut vse funkcije iz $\Ssys_{n\rho}$, a $c_1,\ldots,c_k$ proběgajut vse elementy iz $\Omega$.

Zaisto, dodavanje racionalnyh kombinacij funkcij sistema ne měnja algebraičnogo ranga sistema. $\Lsys_{n\rho}$ ispolnja uslovja 1 i 2, zato jest linearna familija. Dalje, vsaka linearna familija, ktora soderži $\Ssys_{n\rho}$, po 1 i 2 soderži takože $\Lsys_{n\rho}$; zato $\Lsys_{n\rho}$ jest najmenša linearna familija, soderžeča $\Ssys_{n\rho}$.

b) Odpovědajuče, iz linearnoj familije $\Lsys_{n\rho}$ nastava najmenša soderžeča integralna oblast $\Jdom_{n\rho}$ slědujućim potrěbovanjem:

$\Jdom_{n\rho}$ soderži, kromě $\Lsys_{n\rho}$, vse funkcije $h(x)$ i samo te, ktore možno predstaviti kako {\itshape cělu racionalnu kombinaciju} --- s koeficientami iz $\Omega$ --- konečnogo čisla funkcij iz $\Lsys_{n\rho}$, $f_1(x),\ldots,f_k(x)$, pričem $f_1(x),\ldots,f_k(x)$ proběgajut vse funkcije iz $\Lsys_{n\rho}$.

c) Nakonec, iz integralnoj oblasti $\Jdom_{n\rho}$ nastava najmenše soderžeče polje $\Kfield_{n\rho}$ slědujućim potrěbovanjem: $\Kfield_{n\rho}$ soderži, kromě $\Jdom_{n\rho}$, vse funkcije $h(x)$ i samo te, ktore možno predstaviti kako kvocienty dvěh funkcij iz $\Jdom_{n\rho}$.

Ako sobere se tri razširenja v jedinom kroku, dobivamo:

{\itshape Vsaky sistem $\Ssys_{n\rho}$ možno razširiti do najmenšego soderžečego polja $\Kfield_{n\rho}$ slědujućim potrěbovanjem:}

{\itshape $\Kfield_{n\rho}$ soderži, kromě $\Ssys_{n\rho}$, vse funkcije $h(x)$ i samo te, ktore možno predstaviti kako racionalne kombinacije --- s koeficientami iz $\Omega$ --- konečnogo čisla funkcij iz $\Ssys_{n\rho}$, $f_1(x),\ldots,f_k(x)$, pričem $f_1(x),\ldots,f_k(x)$ proběgajut vse funkcije iz $\Ssys_{n\rho}$.}

Iz poslednjego fakta neposrědnje slěduje eksistencija racionalnoj bazy za libovolne oblasti:

Nehaj $\Kfield_{n\rho}$ jest najmenše polje, soderžeče $\Ssys_{n\rho}$, i nehaj racionalna baza $\Kfield_{n\rho}$ jest dana funkcijami
\begin{equation}
  h_1(x),\ h_2(x),\ldots,h_\sigma(x).
\tag{1}
\end{equation}
Po definiciji $\Kfield_{n\rho}$ eksistujut relacije
\begin{equation}
\begin{aligned}
 h_1(x)&=r_1\bigl(g_1(x),\ldots,g_\lambda(x)\bigr);\ldots;\\
 h_\sigma(x)&=r_\sigma\bigl(g_\mu(x),\ldots,g_\nu(x)\bigr),
\end{aligned}
\tag{2}
\end{equation}
kde $r_1,\ldots,r_\sigma$ sut racionalne funkcije svojih argumentov, i
\begin{equation}
  g_1(x),\ldots,g_\lambda(x),\ldots,
  g_\mu(x),\ldots,g_\nu(x)
\tag{3}
\end{equation}
vse prinadležat k $\Ssys_{n\rho}$. Zaradi relacij (2), pokraj (1), takože (3) dava racionalnu bazu $\Kfield_{n\rho}$; a zaradi prinadležnosti funkcij (3) k $\Ssys_{n\rho}$ ona jest jednočasno racionalna baza $\Ssys_{n\rho}$. Zato imajemo:

{\itshape Teorema IV: Libovolny sistem $\Ssys_{n\rho}$ racionalnyh (ili cělyh racionalnyh) funkcij imaje racionalnu bazu algebraičnogo ranga $\rho$, sostavljenu iz konečnogo čisla funkcij samogo sistema.}

Nehaj teper specialno $\Lsys_{n\rho}$ jest linearna familija i
\begin{equation}
  g_1(x),\ldots,g_\rho(x),\ g_{\rho+1}(x),\ldots,g_\nu(x)
\tag{4}
\end{equation}
jest racionalna baza $\Lsys_{n\rho}$. Nehaj naprimer
\[
  g_1(x),\ldots,g_\rho(x)
\]
sut algebraično nezavisne. Ako togda položimo
\[
  g(x)=c_1g_{\rho+1}(x)+c_2g_{\rho+2}(x)+\cdots+c_{\nu-\rho}g_\nu(x),
\]
to $c_i$ možno tako izbrati kako elementy iz $\Omega$, že
\begin{equation}
  g_1(x),\ldots,g_\rho(x),\ g(x)
\tag{5}
\end{equation}
stane racionalnoj bazoj najmenšego soderžečego polja $\Kfield_{n\rho}$. No poněže $\Lsys_{n\rho}$ jest linearna familija, $g(x)$ prinadleži k $\Lsys_{n\rho}$, i (5) dava racionalnu bazu $\Lsys_{n\rho}$; to jest:

{\itshape Teorema V: Racionalnu bazu vsakoj linearnoj familije $\Lsys_{n\rho}$ racionalnyh (ili cělyh racionalnyh) funkcij, a zato osoblivo racionalnu bazu vsakoj integralnoj oblasti $\Jdom_{n\rho}$ i vsakogo polja $\Kfield_{n\rho}$, možno specialno izbrati tako, že ona soderži najviše $\rho+1$ funkcij (i najmanje $\rho$).}

Ograničenje čisla funkcij racionalnoj bazy, najdeno v \S\ 4 za polje, opiraje se zato na svojstvo polja byti linearnoj familijeju; togda kako ograničenje do $\rho$ funkcij v slučaju eksistencije minimalnoj bazy upotrěblja vse predposylki polja.

Treba još zamětiti, že po \S\ 3 (4) vse harakternye svojstva sistemov i najmenših soderžečih sistemov pri odobraženju ostavajut sohranjene.

\end{document}
